\documentclass[UTF8,11pt,a4paper]{ctexart}
\usepackage[a4paper,margin=2.2cm]{geometry}
\usepackage{amsmath,amssymb,mathtools,booktabs,longtable,array}
\usepackage{xcolor,enumitem,hyperref}
\hypersetup{colorlinks=true,linkcolor=blue,urlcolor=blue}
\setlist{itemsep=2pt,topsep=4pt}
\newcommand{\E}{\mathbb E}
\newcommand{\R}{\mathcal R}
\newcommand{\F}{\mathcal F}
\newcommand{\LH}{(L,H)}
\title{曦源项目跨模态理论\\内部小组会讨论稿}
\author{研究交接与接口审查版}
\date{2026 年 9 月}

\begin{document}
\maketitle
\begin{center}
\fbox{\parbox{0.92\textwidth}{\small
这是一份内部讨论稿，不是论文定稿。文中的公式按当前 handoff 材料独立复核；未联网核验文献新颖性，也没有把未定义的项目术语补成既定事实。结论分为：已证明、条件性推导、设计建议和待核验问题。}}
\end{center}

\section{先把四件事连起来}

项目真正要回答的问题不是“高频数据看起来更细，所以一定更有用”，而是：在低频信息已经给定的情况下，高频信息是否还包含可预测目标的增量信息；如果包含，Kronos 和高频分支怎样在不泄漏未来的前提下利用它；如果不包含，增加跨模态结构只会增加成本和过拟合风险。

可以把项目暂时画成一条可检查的链：
\[
 \text{市场观测}\ \longrightarrow\ (L,H)\ \longrightarrow\ (Z_L,Z_H)\ \longrightarrow\ \widehat Y\ \longrightarrow\ \text{预测或交易决策}.
\]
这里 $L$ 是预测时点可用的低频信息（例如 K 线序列），$H$ 是同一时点可用的高频信息（例如逐笔成交或盘口统计），$Z_L=E_L(L)$ 可由 Kronos 产生，$Z_H=E_H(H)$ 是高频编码器输出，最后由预测头 $f$ 产生 $\widehat Y=f(Z_L,Z_H)$。

这条链对应导师和项目的四个接口：
\begin{enumerate}
  \item \textbf{科学问题：} 高频是否有条件增量信息？
  \item \textbf{数学问题：} 增量信息怎样定义、证明或反驳？
  \item \textbf{工程问题：} Kronos 的输入、尺度、预测头和部署时间戳是否满足数学对象的定义？
  \item \textbf{实验问题：} 用什么基线、切分和统计检验判断增量是真实的？
\end{enumerate}

\section{目前最可靠的数学结果}

\subsection{结果 C-CM-001：平方损失下的精确风险恒等式}
\textbf{类型：推导结果；状态：INTERNALLY\_CHECKED；适用条件：有限二阶矩、目标为平方损失。}

定义只用低频信息的最优风险和同时使用低频、高频信息的最优风险：
\[
 \R_L^*=\inf_g \E[(Y-g(L))^2],\qquad
 \R_{L,H}^*=\inf_f \E[(Y-f(L,H))^2].
\]
平方可积条件下，最优预测器分别是条件均值
\[
 g^*(L)=\E[Y\mid L],\qquad f^*(L,H)=\E[Y\mid L,H].
\]
因此有精确恒等式
\[
 \boxed{\R_L^*-\R_{L,H}^*
 =\E\!\left[\big(\E[Y\mid L,H]-\E[Y\mid L]\big)^2\right]\ge 0.}
\]

\paragraph{证明。} 令 $m_H=\E[Y\mid L,H]$、$m_L=\E[Y\mid L]$。条件方差分解给出
\[
 \E[(Y-m_L)^2\mid L]
 =\E[(Y-m_H)^2\mid L]+\E[(m_H-m_L)^2\mid L].
\]
两边取期望即得恒等式。等号成立当且仅当 $m_H=m_L$ 几乎处处，也就是高频信息对平方损失没有条件增量价值。

\paragraph{项目含义。} 这条结果给出“高频有用”的理想判据，但它不保证有限样本、有限模型和有交易成本时一定获益。神经网络只能近似两个条件均值，还必须处理估计误差、优化误差、时间泄漏和成本。

\subsection{结果 C-CM-004：模型 D 的四个闭式量}
\textbf{类型：模型内推导；状态：INTERNALLY\_CHECKED；不是一般金融定理。}

在 handoff 的模型 D 参数化下，$\theta$ 表示高频私有信号比例，$\kappa$ 表示目标对两类潜在成分的载荷参数，$\nu$ 表示噪声方差，并令 $D=2(\theta+\nu)$。四个闭式量为
\[
 u^2=\frac{\bigl(\sqrt{1-\theta}+\kappa\sqrt\theta\bigr)^2}{1+2\kappa^2},\qquad
 R_H=1-\frac{u^2}{1+\nu},
\]
\[
 R=1-\frac{2u^2}{2+\nu-\theta},\qquad
 \Delta I=\frac12\log\!\left(\frac{R_H}{R}\right).
\]
参数范围和原子随机变量的定义必须以模型 D 原始构造为准；这里不能把 $\nu$ 当作 $\kappa$ 代入 $u^2$。在 $\kappa=0$ 时，$u^2=1-\theta$；在 $\kappa=1$ 时，$u^2=(1+2\sqrt{\theta(1-\theta)})/3$。这里 $R_H$ 是单路风险，$R$ 是双路风险；二者不可混用。只有在模型 D 的高斯、最优线性头和单位基准等条件下，$\Delta I$ 才能解释为条件信息增益的对数形式。它不能直接被解释成任意神经表征的“互信息”，也不能把 $\theta$ 自动等同于 $\operatorname{Var}(Z)-1$ 或神经网络通道比例。

\section{三种工程路线，必须分开讨论}

\begin{longtable}{p{0.16\textwidth}p{0.28\textwidth}p{0.26\textwidth}p{0.2\textwidth}}
\toprule
路线 & 结构 & 理论对应 & 组会要确认的事实\\
\midrule
A：输入级融合 & 将高频聚合统计拼入 K 线或 Kronos 输入 & 对象仍是单路 $L'$；当前表征对齐理论不直接适用 & 聚合窗口、事件时间、尺度归一化、是否改变 Kronos 预训练输入分布\\
B：表征级双分支 & Kronos 输出 $Z_L$，高频编码器输出 $Z_H$，预测头同时使用两者 & 最接近 C-CM-001 和 C-CM-004；需要检查配对、尺度、语义与部署可得性 & $Z_L,Z_H$ 是否同一时点配对；是否保留高频私有信息；对齐是否只作用于共享子空间\\
C：输出级辅助任务 & 主预测之外做风险、重构、蒸馏或状态分类 & 这是多任务正则或监督约束，不应自动称为表征对齐 & 辅助任务是否改变目标风险；权重如何预注册；是否引入未来标签或代理泄漏\\
\bottomrule
\end{longtable}

\section{关键约束 X-10：L\textsubscript{align} 到底在对齐什么}

\textbf{当前状态：FORMALIZED 部分完成，尚无普适证明。} “对齐”至少有四种不同含义：边缘分布对齐、条件分布对齐、联合分布对齐、配对样本的表征距离最小化。它们不是同一目标。

若采用配对平方损失，候选形式是
\[
 L_{\rm align}=\E\left[\left\|\phi_H(Z_H)-\phi_L(Z_L)\right\|^2\right],
\]
其中 $\phi_H,\phi_L$ 把两路表征投影到拟共享的子空间。这个形式只有在以下条件大致成立时才有解释：两路样本按同一预测时点配对；投影后的坐标具有可比较尺度；训练时只使用部署时可得的信息；共享子空间确实与目标相关。

最小反例是把高频表征分成共享部分和私有部分：$Z_H=(S_H,P_H)$。若
\[
 \E[Y\mid L,P_H]\ne \E[Y\mid L],
\]
则 $P_H$ 对目标仍有条件增量信息。强行让完整 $Z_H$ 与 $Z_L$ 相同，可能消除这个增量。于是当前最稳妥的设计建议是：只对齐 $\phi_H(Z_H)$ 与 $\phi_L(Z_L)$，把高频私有分量保留给预测头；并用消融实验比较“全量对齐、共享子空间对齐、不对齐”。这是\textbf{设计建议}，不是“对齐必然有效”的定理。

\section{四条理论发展路线的分类}

\begin{enumerate}
  \item \textbf{增量信息路线（最成熟）：} 以 C-CM-001 为基准，估计或检验 $\E[Y\mid L,H]-\E[Y\mid L]$ 是否为零。可转化为嵌套模型比较、时间外推和状态分层评估。
  \item \textbf{接口兼容路线（当前最关键）：} 形式化 X-10，逐项核对 Kronos 输入的时间戳、频率、归一化、预测头和部署可得性。任何一项不满足，C-CM-001 仍是信息论上限，但不能宣称当前实现达到它。
  \item \textbf{对齐安全路线（证明缺口）：} 在共享/私有分解下，研究何种投影对齐不会损害私有增量信息。目前有反例和条件性设计，尚未形成一般性无损定理。
  \item \textbf{决策有效路线（尚待数据）：} 把风险下降转换成成本后的效用：预测误差下降是否超过换手、价差、冲击和执行不确定性的增加。该路线不能由模型 D 的 $\Delta I$ 单独推出。
\end{enumerate}

\section{组会应达成的决定}

请逐项回答并记录证据，而不是用“感觉上可行”代替接口核查：
\begin{enumerate}
  \item 预测时点是什么？高频数据的事件时间、接收时间和实际可用时间分别是什么？
  \item Kronos 的输入究竟是原始 K 线、重采样序列还是已经混入未来窗口的统计量？
  \item 高频分支输出的 $Z_H$ 是否与 $Z_L$ 一一配对？两路尺度和维度如何定义？
  \item 我们选择路线 A、B 还是 C？选择依据是风险、算力、部署延迟还是导师指定接口？
  \item 主结果采用什么基线、损失、时间切分和成本假设？
  \item 哪些话现在可以对导师说，哪些必须保留为“条件性结论”或“待核验”？
\end{enumerate}

建议把最终决定写成一页表格：\textbf{决定、理由、所需数据、否证证据、负责人、截止时间}。如果 X-10 尚未回答，暂不新建跨模态玩具模型；先做接口审计和最小可复现实验。

\section{可以直接对导师说的 60 秒版本}

我们先把问题限定为：在同一个预测时点，给定低频信息 $L$ 后，高频信息 $H$ 是否仍能降低目标的条件均方误差。平方损失下，低频最优风险与双路最优风险之差等于 $\E[(\E[Y\mid L,H]-\E[Y\mid L])^2]$，所以“高频有用”有一个清楚的理想判据。这个判据不等于当前神经网络已经实现了增益：我们还要核对 Kronos 的输入和预测头、两路表征的配对与尺度、实际部署时点，以及交易成本。对齐损失只在共享子空间、时间配对和无泄漏条件下有意义；高频私有信息可能被过度对齐抹掉。下一步先完成 X-10 接口审计，再决定是做双分支表征实验，还是把诊断结果保留为基准。

\section{当前不能声称的内容}

\begin{itemize}
  \item 不能把“扫描中找到一个极小点”写成全局唯一性证明。
  \item 不能把 $R_H$ 和 $R$ 混成同一个风险，也不能把脚本 PASS 当成公式正确的充分证据。
  \item 不能把模型 D 的闭式量直接推广到任意金融数据、任意神经网络或任意交易效用。
  \item 不能在未定位来源时补写文献；“当前材料未找到”只写作“来源未定位”。
  \item 不能把边缘分布对齐、表征距离最小化和预测充分性混为同一个 $L_{\rm align}$。
\end{itemize}

\section{会后最小行动}

优先完成一张 X-10 核查表和一个严格低频基线：固定预测时点，分别训练只用 $L$、只用 $H$、同时用 $(L,H)$ 的模型，使用相同时间切分和调参预算；报告均值、区间、状态分层结果和成本后效用。只有当接口条件明确、双路增益在样本外稳定且不能由泄漏或调参差异解释时，才把跨模态结构升级为项目主线。

\end{document}
